\subsection{Repository-local windows}
\label{sec:exec-polyrepo}

\spec{27} plans one package graph whose packages are owned by several linked
Git repositories. It changes where history and release records are read and
nothing else: the grammar, the bump lattice, the windows, holds,
cancellation, corrections, channel rules, the dependency graph and the
partial-publication rules of \cref{sec:protocol} are unchanged. This
subsection defines the objects the profile adds and states what it promises;
it does not extend the single-history theorems of \cref{sec:guarantees} to
several histories.

\begin{definition}[fleet snapshot and commit identity, \spec{27.2}]
\label{def:exec-fleet}
Before parsing history the engine resolves one immutable map from repository
identity to full commit object identifier, the fixed fleet snapshot; every
participating repository is an initialised, non-shallow checkout at its
pinned commit, and a missing, shallow, moved or ambiguous participant is
\diag{E330}. A commit identity is a pair (repository, full object
identifier); two equal identifier byte strings in different repositories are
incomparable. Commit dates, author dates, modification times, tag creation
dates and fetch order create no ancestry or precedence. The engine reads each
repository's reachable DAG once for the snapshot, preserves native parent
order, and manufactures no synthetic fleet DAG or parent edge. Ancestry,
changed paths, direct pending windows, cancellation barriers, corrections and
\code{Reverts} are evaluated in the repository that owns the commit.
\end{definition}

\begin{definition}[repository-local windows, \spec{27.5}, \spec{27.6}, \spec{27.11}]
\label{def:exec-windows}
For a source repository $r$ with fixed head $h_r$, write $\exRch{r}{x}$ for
the ancestors-or-self of a commit $x$ of $r$, and
\[
  \exafter{r}{x} \;=\; \exRch{r}{h_r} \setminus \exRch{r}{x},
  \qquad
  \exafter{r}{\NONE} \;=\; \exRch{r}{h_r}.
\]
For a consumer package $D$ and a repository $r$, let $s_{D,r}$ be the
revision of $r$ that $D$'s last stable release incorporated and $b_{D,r}$ the
revision of $r$ that $D$'s newest release of any channel incorporated, each
$\NONE$ where no such release exists. $D$'s train and fresh windows over $r$
are
\[
  \exWin{D}{r} \;=\; \exafter{r}{s_{D,r}},
  \qquad
  \exWinf{D}{r} \;=\; \exafter{r}{s_{D,r}} \cap \exafter{r}{b_{D,r}}.
\]
When $r$ owns $D$, the boundaries are the commits of $D$'s own release tags,
and the windows are those of \cref{def:window}. When $r$ does not own $D$,
the boundaries are the positions $D$'s releases prove in $r$, and only these
prove them: a control release checkpoint whose release-commit message names
$D$'s exact tag, changes $D$'s gitlink to the commit peeled from that tag, and
pins the revision of each other repository the release incorporated
(\spec{27.6}); an explicit \code{repositoryBaselines} tuple $(D,
\text{release tag}, r, \text{revision})$, with stable and prerelease tags as
separate entries (\spec{27.6}); or, in the linked peer topology, the route of
fleet links recorded hop by hop from the tree of the release commit that
names $D$'s tag (\spec{27.11}). A matching object identifier alone proves
nothing. Missing or ambiguous evidence is \diag{E333}, never a guess and
never a choice by date or proximity. Admission of a unit $u$ carried by a
commit of $r$ at a consumer $D$ remains $\cof{u} \in \exWinf{D}{r}$: $D$'s
last release had not yet incorporated the qualified revision. A source unit
resolves directly only packages of its own repository and propagates through
the one combined graph; delivery (\spec{13.4a}) compares $D$'s boundary and
the release commits of the provider in the provider's repository. A window
over $r$ is a function of $r$, $h_r$ and one boundary revision of $r$ and of
nothing else, so packages of any owners share an immutable train window over
$r$ exactly when their stable boundary there is the same qualified revision,
and a fresh window exactly when both boundaries are.
\end{definition}

The identity $\exWinf{D}{r} = \exafter{r}{s_{D,r}} \cap \exafter{r}{b_{D,r}}$
is the fresh-window identity of \cref{prop:marker-pass}, read per
repository; it is what lets \cref{sec:complexity} key every window by one
boundary revision. Newest-wins precedence between competing directives is
defined within one repository's DAG and within the control history; source
revisions of different repositories have no such order, and where the
existing rule requires a single winner among incomparable revisions the
engine fails with \diag{E334} rather than choosing by date, traversal order,
repository name or identifier (\spec{27.4}). An explicit control directive can
resolve the choice, evaluated against the gitlinks its own commit pinned, its
causal snapshot, so that it cannot rewrite work the control revision had not
yet observed; in the linked peer topology no repository addresses the whole
fleet and the conflict is refused outright.

\begin{definition}[recording order, \spec{27.7}, \spec{27.11}]
\label{def:exec-recording}
The combined publish sequence is that of \cref{def:order} across all
repositories, and publication remains non-atomic. Each successful package
receives its immutable release tag in its owning repository at the ordinary
tag-after-publish point of \cref{def:tagging}. Where a source release commit
and push are enabled, the source release is recorded durably first, and only
after its branch and tags are reachable from the configured remote may a
control gitlink or a peer's pin advance to it: a pushed pointer never names a
commit nobody else can fetch. A control checkpoint, when enabled and
non-empty, follows the source record and never precedes it; no empty commit
is forced and no checkpoint is required. A failure to record or durably push
a tag is \diag{E335}; prior successes remain successful and dependents are
protected by \cref{def:blocking}. If the source record succeeds and the
control checkpoint then fails, \diag{E335} names the repository, the full
revision and the tag and requires explicit repair; the next run remains
\diag{E330} while the checkout and the committed gitlink are unpinned, and
the engine neither auto-commits the repair, nor forces a checkpoint, nor
republishes the source, nor infers success from the worktree. A successful
source tag is never deleted, moved or duplicated to make a checkpoint look
atomic, and no rollback is inferred. In the linked peer topology the route
evidence is settled by ordinary commits before publication, deepest hop
first, all or nothing per package (\diag{E333} otherwise), and each
settlement is an admitted transition of the repository it was written in.
The fleet lock covers the control repository, every active source and every
standalone package as one fleet (\diag{E336} otherwise), is taken in one
stable total order over lock identities, and is released in reverse.
\end{definition}

\paragraph{Bounds.}
\spec{27.9} requires the engine to index the control gitlink transitions
once per fixed control snapshot, not once per consumer or provider pair; to
store each parsed (repository, commit) record once; to share window
membership only under the keys of \cref{def:exec-windows}; and to bound its
reachability and ancestry caches, because a cache of every queried pair can
grow quadratically in one repository's history even when that history is
walked once. With the notation of \cref{sec:complexity}, input traversal is
$O(G + \sum_q (H_q + A_q))$, window reachability is $O(\sum_q (H_q + A_q)\cdot
\mathrm{bw}(m_q) + I_w)$ with one marker pass per repository, and in the linked
peer topology boundary resolution is $O(Q + X \cdot Y)$ lookups over $B$ tree
reads, with settlement in at most $Q - 1$ commits and at most two in a star.
Publication revalidation adds $O(\sum_q J_q (1 + T_q))$ comparison work, with
$J_q$ the fleet or package checks repository $q$ takes part in and $T_q$ its
relevant tag refs, and retaining every package's input closure costs memory
linear in the total number of repository memberships across those closures,
at most $P \cdot Q$. The profile makes no globally linear end-to-end claim and permits
no synthetic ancestry shortcut.

\begin{remark}[what this subsection does not prove]
\label{rem:exec-multihistory}
\Cref{lem:window}, and every theorem of \cref{sec:guarantees} that rests on
it, is stated for one history with one \Head{}. Under \cref{def:exec-windows}
a consumer has one window per repository its plan reads, each measured from a
boundary its own release proves there, and a release of the consumer moves
every one of those boundaries at once only when the evidence of
\cref{def:exec-recording} is written. Restating the retry invariant over the
fleet snapshot and proving the eight guarantees per repository is a separate
piece of work; this report defines the objects, and \spec{27} promises that
the single-history rules apply unchanged inside each repository.
\end{remark}
